
\chapter{Problem Statement}
\label{chap:ps}
\chaptermark{Second Chapter Heading}

\section{Background of the Problem}
In this section, provide a brief background to the problem your research addresses. Describe the broader context and why this issue is significant. For example, you could explain the challenges faced by industries in robotics, automation, or welding, and why they are relevant to your study. Discuss any existing limitations or gaps in current solutions.

Clearly define the specific problem your research aims to solve. For example, you might identify inefficiencies in current welding processes, limitations of existing robotic systems, or challenges with integrating vision-based systems into industrial workflows.

\section{Impact of the Problem}
Explain the impact of this problem on the company, industry, or specific field of study. Discuss how the problem affects productivity, quality, safety, or costs. For example, you could highlight how manual welding is time-consuming and prone to errors, leading to increased operational costs and lower quality.

\section{Scope of the Problem}
Narrow down the scope of the problem to focus on what your project addresses. Explain the specific aspect of the problem that you will tackle in your research. For instance, your project may focus on automating a specific part of the welding process or improving the accuracy of robotic vision systems for better alignment.

\section{Summary of the Problem}
Summarize the problem in one or two sentences to clearly state what your research aims to solve. This will set the stage for the objectives and methodology you will describe later in your thesis.



